%%=====================================================================
%% SBESC 2026 -- XVI Symposium on Computing Systems Engineering
%% VERSAO EM PORTUGUES (papers em PT sao publicados na SBC OpenLib).
%% Template: IEEE two-column conference (maximo 6 paginas).
%% Submissao: JEMS3 -> https://jems3.sbc.org.br/events/628 (prazo: 2026-07-10)
%%
%% COMO GERAR O PDF:
%%   Recomendado: subir este .tex no Overleaf, compilador pdfLaTeX,
%%   clicar Recompile e baixar o PDF. Nao requer imagens externas
%%   (a figura da arquitetura e desenhada com TikZ).
%%
%% ANTES DE SUBMETER (conferir):
%%   - e-mails dos autores (o segundo e um placeholder)
%%   - respeitar o limite de 6 paginas no PDF final
%%=====================================================================
\documentclass[conference]{IEEEtran}
\IEEEoverridecommandlockouts

\usepackage[utf8]{inputenc}
\usepackage[T1]{fontenc}
\usepackage[brazil]{babel}
\usepackage{cite}
\usepackage{amsmath,amssymb,amsfonts}
\usepackage{graphicx}
\usepackage{booktabs}
\usepackage{array}
\usepackage{tikz}
\usetikzlibrary{positioning,arrows.meta}
\usepackage[hidelinks]{hyperref}
\usepackage{url}

\begin{document}

\title{Autenticação Contextual Autônoma para Aplicações\\ Web Combinando um Ensemble de IA e\\ Auditoria em Blockchain Permissionada}

\author{\IEEEauthorblockN{João Guedes de Brito}
\IEEEauthorblockA{\textit{Programa de Pós-Graduação em Engenharia de Sistemas e Produtos (PPGESP)}\\
\textit{Instituto Federal da Bahia (IFBA)}\\
Salvador, Bahia, Brasil \\
joaoguedesdebrito@gmail.com}
\and
\IEEEauthorblockN{Cleber Jorge Lira de Santana}
\IEEEauthorblockA{\textit{Programa de Pós-Graduação em Engenharia de Sistemas e Produtos (PPGESP)}\\
\textit{Instituto Federal da Bahia (IFBA)}\\
Salvador, Bahia, Brasil \\
cleber.santana@ifba.edu.br}
}

\maketitle

\begin{abstract}
Mecanismos tradicionais de autenticação baseados em senhas estáticas ou
em autenticação multifator (MFA) fixa permanecem amplamente adotados,
mas ignoram o contexto comportamental e ambiental de cada tentativa de
acesso e são cada vez mais explorados por ataques de \emph{credential
stuffing}, phishing e tomada de conta. Este artigo apresenta um sistema
autônomo de autenticação contextual para aplicações web que combina um
ensemble de aprendizado de máquina com auditoria em blockchain
permissionada. Cada login é descrito por quatorze critérios contextuais
e comportamentais, avaliados por um ensemble de Isolation Forest, Random
Forest e uma rede neural multicamada; o escore de risco resultante
conduz uma decisão adaptativa em três faixas (permitir, reforçar com MFA
ou bloquear). Cada decisão é registrada de forma imutável em uma rede
Hyperledger Fabric, gerando uma trilha de auditoria verificável e de não
repúdio, cuja necessidade frente a logs assinados e bancos append-only é
explicitamente justificada. Em avaliação controlada sobre um conjunto de
dados sintéticos com cerca de 50.000 registros de login (95\% legítimos,
5\% suspeitos), o ensemble alcançou 98,1\% de acurácia, 80,6\% de
F1-score e AUC-ROC de 0,99 no conjunto de teste, com taxa de falsos
positivos de 1,0\% e latência de inferência de aproximadamente 18,6\,ms
por requisição. Esses resultados são preliminares e baseiam-se em dados
sintéticos; caracterizam a viabilidade técnica da abordagem, e não um
estudo de campo, devendo ser confirmados com dados reais.
\end{abstract}

\begin{IEEEkeywords}
autenticação contextual, autenticação baseada em risco, aprendizado de
máquina, ensemble, detecção de anomalias, blockchain, trilha de
auditoria, segurança da informação.
\end{IEEEkeywords}

\section{Introdução}
A autenticação é a primeira linha de defesa de aplicações web.
Historicamente, apoiou-se em um fator de conhecimento --- a senha ---,
cujas fragilidades são amplamente documentadas: senhas são reutilizadas
entre serviços, vazadas em larga escala e exploradas por força bruta,
\emph{credential stuffing} e phishing~\cite{nist80063b,verizon2024}. A
autenticação multifator (MFA) mitiga parte dessa exposição ao combinar
fatores de conhecimento, posse e inerência, mas a MFA estática ainda
sucumbe a engenharia social, troca de chip (SIM swapping) e interceptação
de códigos em tempo real, além de introduzir atrito na experiência do
usuário~\cite{nist80063b}.

Para enfrentar essas limitações, as diretrizes de identidade digital do
NIST e o modelo de Zero Trust recomendam mecanismos \emph{adaptativos},
nos quais o nível de verificação exigido é proporcional ao risco
estimado de cada tentativa~\cite{nist80063b,rose2020}. Isso desloca a
autenticação de uma decisão binária e estática para uma avaliação
contínua e sensível ao contexto --- abordagem conhecida na literatura
como autenticação baseada em risco (\emph{Risk-Based Authentication},
RBA). Estudos empíricos de serviços em produção mostram, contudo, que o
RBA implantado é dominado por um único sinal (tipicamente o endereço IP)
e por heurísticas ajustadas manualmente~\cite{wiefling2019,wiefling2023}.
Em paralelo, quando as decisões de acesso são delegadas a modelos
autônomos, elas precisam ser \emph{auditáveis} por terceiros (por
exemplo, em fiscalizações de conformidade), requisito que registros de
decisão sob a guarda de um único operador não atendem plenamente.

Este artigo propõe um sistema autônomo de autenticação contextual que
unifica, em um único fluxo, uma decisão adaptativa por aprendizado de
máquina (ML) e sua auditoria imutável. As principais contribuições são:
\begin{itemize}
  \item uma arquitetura ponta a ponta que funde quatorze critérios
        contextuais e comportamentais em um único escore de risco
        calculado por um ensemble de Isolation Forest, Random Forest e
        rede neural, conduzindo uma decisão adaptativa em três faixas;
  \item uma camada de auditoria em blockchain permissionada acoplada à
        decisão, com justificativa explícita do \emph{porquê} da
        necessidade de blockchain frente a logs assinados ou bancos
        append-only; e
  \item uma metodologia de avaliação reprodutível, com baselines e um
        estudo preliminar de viabilidade sobre dados sintéticos.
\end{itemize}

O restante do artigo organiza-se assim. A Seção~\ref{sec:related} revisa
os trabalhos correlatos. A Seção~\ref{sec:arch} descreve a arquitetura
proposta; a Seção~\ref{sec:ai}, o motor de IA; e a Seção~\ref{sec:bc}, a
camada de auditoria em blockchain. A Seção~\ref{sec:method} apresenta a
metodologia de avaliação e a Seção~\ref{sec:results}, os resultados. A
Seção~\ref{sec:conclusion} conclui.

\section{Trabalhos Correlatos}
\label{sec:related}
\textbf{Biometria comportamental.} A autenticação contínua a partir de
padrões de interação está bem estabelecida. O TypeNet~\cite{acien2022} e
o TypeFormer~\cite{stragapede2024} identificam usuários pela dinâmica de
digitação com baixas taxas de erro, e o Touchalytics~\cite{frank2013}
mostra que padrões de toque discriminam usuários em dispositivos móveis.
Esses trabalhos confirmam o poder discriminativo do comportamento do
usuário, mas concentram-se em uma única modalidade biométrica e não
tratam da auditoria imutável das decisões.

\textbf{Autenticação baseada em risco.} Wiefling \emph{et
al.}~\cite{wiefling2019} conduziram o primeiro estudo empírico de larga
escala sobre RBA em serviços reais, e trabalhos posteriores
\cite{wiefling2023} mostraram que os fatores de risco devem ser
cuidadosamente ajustados por serviço. Nossa proposta segue essa linha,
mas substitui heurísticas fixas por um ensemble de ML e adiciona uma
camada de não repúdio.

\textbf{Detecção de anomalias.} O Isolation Forest~\cite{liu2008} é
referência para detecção não supervisionada de outliers em alta
dimensionalidade, e modelos supervisionados e profundos
\cite{goodfellow2016} o complementam quando há rótulos; Killourhy e
Maxion~\cite{killourhy2009} estabeleceram um marco metodológico para a
detecção de anomalias em dinâmica de digitação. Combinamos ambas as
naturezas --- não supervisionada e supervisionada --- em um ensemble.

\textbf{Blockchain para auditoria.} Embora a auditoria de logs baseada em
blockchain venha sendo investigada, a maioria das abordagens trata o
registro de eventos de forma \emph{desacoplada} do mecanismo de decisão.
O diferencial deste trabalho é integrar, em um único fluxo, a decisão
adaptativa por IA e o seu registro imutável em rede permissionada. A
Tabela~\ref{tab:related} sintetiza trabalhos representativos: nenhuma
linha de pesquisa reúne simultaneamente (i)~a fusão de múltiplos sinais
contextuais e comportamentais, (ii)~a decisão adaptativa por ensemble e
(iii)~a auditoria imutável acoplada à decisão.

\begin{table}[t]
\caption{Posicionamento Frente a Trabalhos Correlatos Representativos}
\label{tab:related}
\centering
\renewcommand{\arraystretch}{1.2}
\begin{tabular}{@{}p{1.6cm}p{2.0cm}p{2.95cm}@{}}
\toprule
\textbf{Trabalho} & \textbf{Técnica} & \textbf{Limitação principal} \\
\midrule
TypeNet~\cite{acien2022} & Rede neural recorrente & Modalidade única; sem auditoria da decisão \\
Touchalytics~\cite{frank2013} & SVM / $k$-NN & Restrito a telas de toque \\
RBA~\cite{wiefling2019} & Heurísticas/estatística & Regras fixas; sem IA adaptativa \\
Isolation Forest~\cite{liu2008} & Não supervisionado & Sinal isolado; não decide o acesso \\
\textbf{Este trabalho} & Ensemble (IF+RF+DL) + blockchain & Integra os três eixos em um fluxo \\
\bottomrule
\end{tabular}
\end{table}

\section{Arquitetura Proposta}
\label{sec:arch}
A arquitetura articula microsserviços, gestão federada de identidade,
análise de risco por IA e registro imutável em blockchain, seguindo os
princípios de segurança por desenho e verificação contínua do Zero
Trust~\cite{rose2020,nist80053r5}. Em vez de uma decisão binária no
momento do login, o modelo adota uma perspectiva contínua e orientada a
risco. Organiza-se em cinco camadas integradas (Fig.~\ref{fig:arch}):

\begin{enumerate}
  \item \textbf{Aplicação Modular.} Serviços independentes (Spring Boot)
        e front end Angular expondo APIs REST autenticadas.
  \item \textbf{Identidade e Acesso.} Autenticação federada via Keycloak
        (OAuth2 / OpenID Connect), emissão de JWT e suporte a MFA.
  \item \textbf{Inteligência Contextual.} Extração dos quatorze critérios
        comportamentais e cálculo do escore de risco pelo ensemble.
  \item \textbf{Decisão Adaptativa.} Aplicação dos limiares de risco,
        resultando em \emph{permitir}, \emph{reforçar com MFA} ou
        \emph{bloquear}.
  \item \textbf{Registro Imutável e Auditoria.} Persistência de cada
        decisão em blockchain permissionada, garantindo rastreabilidade
        e não repúdio.
\end{enumerate}

O fluxo operacional é uma cadeia decisória encadeada: a tentativa de
login dispara a extração dos critérios contextuais; o motor de IA calcula
o escore de risco; os limiares de decisão são aplicados; o resultado é
registrado na blockchain; e o sistema permite o acesso, exige MFA ou o
bloqueia. A separação em camadas permite evoluir cada responsabilidade de
forma independente e aplicar políticas de segurança específicas em cada
nível.

\begin{figure}[t]
\centering
\begin{tikzpicture}[
  layer/.style={draw, rounded corners, minimum width=7.4cm,
    minimum height=0.72cm, align=center, font=\footnotesize, fill=gray!8},
  arr/.style={-{Latex[length=2mm]}, thick}
]
\node[layer] (l1) {C1 -- Aplicação Modular \\ \scriptsize Serviços Spring Boot $\cdot$ front end Angular $\cdot$ APIs REST};
\node[layer, below=0.28cm of l1] (l2) {C2 -- Identidade e Acesso \\ \scriptsize Keycloak $\cdot$ OAuth2 / OIDC $\cdot$ JWT $\cdot$ MFA};
\node[layer, below=0.28cm of l2] (l3) {C3 -- Inteligência Contextual \\ \scriptsize 14 critérios $\cdot$ ensemble IF + RF + DL $\rightarrow$ escore};
\node[layer, below=0.28cm of l3] (l4) {C4 -- Decisão Adaptativa \\ \scriptsize limiares $\rightarrow$ permitir / MFA / bloquear};
\node[layer, below=0.28cm of l4] (l5) {C5 -- Registro Imutável e Auditoria \\ \scriptsize Hyperledger Fabric $\cdot$ trilha de auditoria verificável};
\draw[arr] (l1) -- (l2);
\draw[arr] (l2) -- (l3);
\draw[arr] (l3) -- (l4);
\draw[arr] (l4) -- (l5);
\end{tikzpicture}
\caption{Arquitetura em cinco camadas do sistema proposto.}
\label{fig:arch}
\end{figure}

\section{Motor de Decisão por IA}
\label{sec:ai}
\textbf{Ensemble.} O motor de risco combina três aprendizes
complementares. O \emph{Isolation Forest}~\cite{liu2008} atua como
primeiro nível de análise, isolando outliers no espaço multidimensional
de critérios com poucas partições aleatórias; contribui com 40\% do peso
final. O \emph{Random Forest} contribui com 30\%, capturando relações não
lineares entre os critérios e mantendo-se interpretável. Uma \emph{rede
neural multicamada}~\cite{goodfellow2016} contribui com os 30\%
restantes, com camada de entrada de 14 neurônios, duas camadas ocultas
(64 e 32 unidades, ativações ReLU, regularização por dropout) e saída
sigmoid. O escore final é a média ponderada
\begin{equation}
S = 0{,}4\,S_{\mathrm{IF}} + 0{,}3\,S_{\mathrm{RF}} + 0{,}3\,S_{\mathrm{DL}},
\label{eq:ensemble}
\end{equation}
em que cada componente retorna um valor em $[0,1]$ (0: normal, 1:
altamente anômalo). A combinação de aprendizes não supervisionados e
supervisionados aumenta a robustez tanto a ataques inéditos (zero-day)
quanto a fraudes conhecidas.

\textbf{Critérios.} Cada tentativa é representada por quatorze critérios
em quatro famílias: \emph{temporais} (hora do acesso em escala circular,
dia da semana, intervalo desde o último login, frequência em 24\,h);
\emph{geográficos/rede} (distância geodésica do local habitual, tipo de
rede, ASN, velocidade de deslocamento impossível);
\emph{dispositivo/navegador} (fingerprint do navegador, sistema
operacional, tipo de dispositivo, detecção de automação/headless); e
\emph{comportamento de sessão} (padrão de digitação e histórico de falhas
recentes). Esses sinais visam credential stuffing, viagem impossível,
sequestro de sessão, bots e desvios zero-day.

\textbf{Limiares adaptativos.} O escore é mapeado em três faixas de
risco, calibradas empiricamente para equilibrar segurança e usabilidade:
baixo risco ($S<0{,}3$) permite o acesso sem etapas adicionais; risco
moderado ($0{,}3\le S<0{,}7$) aciona MFA; e alto risco ($S\ge0{,}7$)
bloqueia a tentativa, registra o evento na blockchain e notifica o
administrador de segurança. Em produção, o motor é exposto como um
microsserviço independente (FastAPI) consumido pelo backend via REST,
mantendo os três modelos em memória para inferência paralela. Um
mecanismo de fallback heurístico preserva a disponibilidade caso o
serviço de IA fique inacessível. Os modelos são reentreinados
periodicamente para acompanhar a mudança de comportamento (concept
drift)~\cite{gama2014}, e todo modelo treinado é versionado para permitir
rollback.

\section{Auditoria Baseada em Blockchain}
\label{sec:bc}
Cada decisão é registrada de forma imutável em uma rede permissionada
baseada em Hyperledger Fabric~\cite{fabric2024}, cujo controle de
identidade e governança estruturada são adequados a ambientes
corporativos e governamentais~\cite{nistir8202}. Para cada evento, o
registro armazena o hash dos critérios de entrada, o escore de risco, a
decisão (permitir, MFA ou bloquear) e o timestamp da transação, encadeado
ao registro anterior de modo a formar uma cadeia cronológica de
evidências.

A escolha do blockchain frente a mecanismos mais simples é justificada, e
não presumida. Logs assinados digitalmente e bancos append-only já
oferecem integridade e não repúdio por meio de assinaturas
criptográficas, sendo suficientes quando existe uma autoridade central
confiável que opera o registro. A diferença essencial de um blockchain
permissionado \emph{não} está na assinatura em si, mas em não depender da
confiança em um único custodiante: o registro é replicado e validado por
consenso entre múltiplos participantes, tornando-se resistente à
adulteração mesmo por administradores privilegiados do próprio sistema e
permitindo verificação independente por auditores externos. Isso é
determinante aqui, pois as decisões de acesso são tomadas de forma
autônoma por modelos de IA e precisam ser auditáveis por terceiros (por
exemplo, em fiscalizações de conformidade da LGPD~\cite{lgpd2018}),
cenário em que o operador, sozinho, não deve poder alterar a trilha de
decisões. Reconhece-se o custo --- maior complexidade operacional e menor
throughput do que um banco append-only ---, que se justifica justamente
em cenários multipartes e regulados, delimitando o escopo em que o
blockchain é efetivamente necessário.

\section{Metodologia de Avaliação}
\label{sec:method}
\textbf{Dados e rotulagem.} Diante da ausência de bases públicas
rotuladas para autenticação contextual, construiu-se um conjunto de dados
sintéticos com cerca de 50.000 registros de login, reproduzindo a
distribuição observada em produção (cerca de 95\% de acessos legítimos e
5\% suspeitos). Cada registro contém os quatorze critérios da
Seção~\ref{sec:ai}. Os rótulos foram derivados de regras heurísticas que
codificam padrões conhecidos de fraude (viagem impossível, dispositivo
desconhecido, rajadas de falhas de login) e revisados por especialistas
em segurança. Por se tratar de dados sintéticos, não há coleta de
informações pessoais identificáveis nesta fase; a evolução para dados
reais seguirá os princípios de minimização e anonimização exigidos pela
LGPD~\cite{lgpd2018}.

\textbf{Protocolo.} O conjunto foi particionado de forma estratificada em
70\% para treinamento e 30\% para teste. Validação cruzada estratificada
com 5 folds foi aplicada sobre a partição de treinamento para seleção de
hiperparâmetros, reservando a partição de teste exclusivamente para a
avaliação final, evitando vazamento de dados (data leakage). O desempenho
foi medido por acurácia, precisão, recall, F1-score, taxa de falsos
positivos (FPR), AUC-ROC e latência de decisão. As sementes aleatórias
foram fixadas, os modelos e hiperparâmetros versionados, e todo o
ambiente foi containerizado (Docker/Kubernetes) para reprodução
determinística.

\textbf{Baselines.} A Tabela~\ref{tab:baseline} confronta a proposta com
baselines representativos da literatura e da prática. Os valores dos
baselines provêm da literatura citada e não foram reexecutados aqui; a
comparação é, portanto, qualitativa quanto às capacidades de cada
abordagem.

\begin{table}[t]
\caption{Comparação com Abordagens de Referência}
\label{tab:baseline}
\centering
\renewcommand{\arraystretch}{1.2}
\begin{tabular}{@{}p{2.05cm}p{1.45cm}p{3.05cm}@{}}
\toprule
\textbf{Abordagem} & \textbf{Contextual?} & \textbf{Limitação principal} \\
\midrule
Senha + 2FA estático & Não & Ignora contexto; phishing/SIM swap \\
RBA por IP~\cite{wiefling2019} & Parcial & Feature única; ajuste manual \\
Keystroke isolado~\cite{killourhy2009} & Sim & Offline, fator único (EER $\approx$9,6\%) \\
\textbf{Proposta (ensemble)} & Sim & Adaptativa, explicável e auditável \\
\bottomrule
\end{tabular}
\end{table}

\section{Resultados e Discussão}
\label{sec:results}
Sobre o conjunto de teste (30\% dos dados, não vistos durante o
treinamento), o ensemble alcançou acurácia de 98,1\%, precisão de 80,7\%
e recall de 80,4\% na detecção de acessos suspeitos, com F1-score de
80,6\% e AUC-ROC de 0,99, indicando forte capacidade de ordenação entre
acessos legítimos e suspeitos. A taxa de falsos positivos manteve-se
baixa, em 1,0\%, fator relevante para a usabilidade. A validação cruzada
estratificada com 5 folds confirmou a estabilidade dos resultados
(acurácia $98,1\%\pm0,1\%$; F1 $80,5\%\pm0,7\%$). A
Tabela~\ref{tab:results} resume esses números.

\begin{table}[t]
\caption{Desempenho do Ensemble no Conjunto de Teste}
\label{tab:results}
\centering
\renewcommand{\arraystretch}{1.2}
\begin{tabular}{@{}lc@{\hspace{5mm}}lc@{}}
\toprule
\textbf{Métrica} & \textbf{Valor} & \textbf{Métrica} & \textbf{Valor} \\
\midrule
Acurácia   & 98,1\% & F1-score        & 80,6\% \\
Precisão   & 80,7\% & AUC-ROC         & 0,99   \\
Recall     & 80,4\% & Taxa falsos pos.\ & 1,0\% \\
\midrule
\multicolumn{2}{@{}l}{Acurácia (5-fold)} & \multicolumn{2}{l}{$98,1\%\pm0,1\%$} \\
\multicolumn{2}{@{}l}{F1-score (5-fold)} & \multicolumn{2}{l}{$80,5\%\pm0,7\%$} \\
\multicolumn{2}{@{}l}{Latência do ensemble} & \multicolumn{2}{l}{$\approx$18,6\,ms/req.} \\
\bottomrule
\end{tabular}
\end{table}

Quanto à cadeia de autenticação, a integração entre o Keycloak e o
serviço de ML foi validada com sucesso. A inferência do ensemble foi
medida em cerca de 18,6\,ms por requisição, muito abaixo do orçamento de
200\,ms definido para o componente de IA; o fluxo completo de
autenticação contextual --- incluindo a validação de identidade e o
registro em blockchain --- permaneceu na ordem de algumas centenas de
milissegundos. O pipeline de registro em blockchain mostrou-se consistente
nos testes de imutabilidade, com validação criptográfica bem-sucedida em
todas as transações registradas, e o ELK Stack e o Grafana permitiram a
visualização em tempo real dos eventos de autenticação.

\textbf{Discussão.} Os resultados parciais são promissores e sugerem a
viabilidade técnica da abordagem; contudo, por decorrerem de dados
sintéticos e de uma fase preliminar, não permitem ainda confirmar a
hipótese central de forma conclusiva. Uma comparação rigorosa com os
baselines e a validação com dados reais permanecem necessárias antes de
qualquer afirmação quantitativa de superioridade. O F1-score de 80,6\%,
embora adequado ao estágio atual, evidencia espaço para melhorar o recall
quando os dados reais substituírem o conjunto sintético. Ainda assim, em
ambiente controlado, a abordagem adaptativa e orientada a contexto
detectou padrões de acesso suspeitos --- como os cenários simulados de
credential stuffing e sequestro de sessão --- que escapam a sistemas
baseados apenas em senha e 2FA estático, enquanto a camada de blockchain
tornou inviável a adulteração retroativa dos registros sem detecção.

\section{Considerações Finais e Trabalhos Futuros}
\label{sec:conclusion}
Este artigo apresentou um sistema autônomo de autenticação contextual que
unifica, em um único fluxo, uma decisão adaptativa por IA e sua auditoria
imutável em blockchain permissionada. As principais contribuições são a
arquitetura ponta a ponta que funde quatorze sinais contextuais e
comportamentais por meio de um ensemble, a camada de auditoria acoplada à
decisão com justificativa explícita para o uso de blockchain e uma
metodologia de avaliação reprodutível. Experimentos preliminares sobre
dados sintéticos indicaram viabilidade técnica (98,1\% de acurácia, 80,6\%
de F1-score, AUC de 0,99, latência de inferência de $\approx$18,6\,ms),
constituindo uma validação metodológica e não um substituto para a
avaliação com dados reais. Como trabalhos futuros, pretende-se substituir
o conjunto sintético por dados reais ou semirreais coletados em ambiente
controlado, refinar a calibração de modelos e limiares, avaliar
desempenho e escalabilidade sob carga realista e aprofundar a análise de
segurança frente ao modelo de ameaças.

\begin{thebibliography}{00}
\bibitem{nist80063b} National Institute of Standards and Technology, ``Digital Identity Guidelines: Authentication and Lifecycle Management,'' NIST SP 800-63B, 2017.
\bibitem{verizon2024} Verizon, ``2024 Data Breach Investigations Report (DBIR),'' Verizon Business, 2024.
\bibitem{rose2020} S. Rose, O. Borchert, S. Mitchell, and S. Connelly, ``Zero Trust Architecture,'' NIST SP 800-207, 2020.
\bibitem{wiefling2019} S. Wiefling, L. Lo Iacono, and M. D\"urmuth, ``Is This Really You? An Empirical Study on Risk-Based Authentication Applied in the Wild,'' in \emph{IFIP SEC}, 2019, pp. 134--148.
\bibitem{wiefling2023} S. Wiefling, P. R. J\o rgensen, S. Thunem, and L. Lo Iacono, ``Pump Up Password Security! Evaluating and Enhancing Risk-Based Authentication on a Real-World Large-Scale Online Service,'' \emph{ACM Trans. Priv. Secur.}, vol. 26, no. 1, 2023.
\bibitem{acien2022} A. Acien, A. Morales, R. Vera-Rodriguez, J. Fierrez, and J. V. Monaco, ``TypeNet: Deep Learning Keystroke Biometrics,'' \emph{IEEE Trans. Biom. Behav. Identity Sci.}, vol. 4, no. 1, pp. 57--70, 2022.
\bibitem{stragapede2024} G. Stragapede \emph{et al.}, ``TypeFormer: Transformers for Mobile Keystroke Biometrics,'' \emph{Pattern Recognition Letters}, 2024.
\bibitem{frank2013} M. Frank, R. Biedert, E. Ma, I. Martinovic, and D. Song, ``Touchalytics: On the Applicability of Touchscreen Input as a Behavioral Biometric for Continuous Authentication,'' \emph{IEEE Trans. Inf. Forensics Security}, vol. 8, no. 1, pp. 136--148, 2013.
\bibitem{liu2008} F. T. Liu, K. M. Ting, and Z.-H. Zhou, ``Isolation Forest,'' in \emph{IEEE Int. Conf. Data Mining (ICDM)}, 2008, pp. 413--422.
\bibitem{killourhy2009} K. S. Killourhy and R. A. Maxion, ``Comparing Anomaly-Detection Algorithms for Keystroke Dynamics,'' in \emph{IEEE/IFIP DSN}, 2009, pp. 125--134.
\bibitem{goodfellow2016} I. Goodfellow, Y. Bengio, and A. Courville, \emph{Deep Learning}. Cambridge, MA: MIT Press, 2016.
\bibitem{gama2014} J. Gama, I. \v{Z}liobait\.e, A. Bifet, M. Pechenizkiy, and A. Bouchachia, ``A Survey on Concept Drift Adaptation,'' \emph{ACM Comput. Surv.}, vol. 46, no. 4, pp. 1--37, 2014.
\bibitem{fabric2024} Hyperledger Foundation, ``Hyperledger Fabric Documentation, v2.5,'' Linux Foundation, 2024.
\bibitem{nistir8202} National Institute of Standards and Technology, ``Blockchain Technology Overview,'' NISTIR 8202, 2018.
\bibitem{nist80053r5} National Institute of Standards and Technology, ``Security and Privacy Controls for Information Systems and Organizations,'' NIST SP 800-53 Rev.\ 5, 2020.
\bibitem{lgpd2018} Brasil, ``Lei n.\ 13.709/2018 -- Lei Geral de Proteção de Dados Pessoais (LGPD),'' 2018.
\end{thebibliography}

\end{document}
